\section{Introduction}
\label{sec:intro}

Systems built from many language-model agents are now routine: agents that
debate, review each other's work, negotiate, and simulate populations
\citep{park2023generative,du2023debate,chuang2024simulating}. Each agent reads
what others produce, forms some judgement of how much to rely on them, and
adjusts its own position. That is a society, and societies have collective states
that no individual chose. The question this paper asks is whether one of the
collective states available to such a system is \emph{discrimination}: a
population that sorts itself into mutually distrustful groups along a label that
carries no information.

The usual explanation for group bias in machine learning locates it in the data:
a model trained on a biased corpus reproduces the bias. That account cannot
explain a population-level effect, because it says nothing about interaction. We
therefore ask a different question. Suppose every agent is individually
well-specified --- learning by an algorithm that is provably optimal in the
setting it was designed for --- and suppose the only defect is that some agents
attend to one irrelevant feature of a message, namely who sent it. What does the
\emph{population} do?

We study this in a model in which the analysis can be carried out: a society of
$N$ agents, each carrying a perceptron over a $K$-dimensional embedding of
issues, interacting by stating opinions about issues drawn from a shared agenda
of size $P$. Two features make the model unusually apt for the question. First,
agents can agree on some issues and disagree on others, so ideological distance
is graded rather than binary. Second, learning is derived rather than posited:
the update follows from a Bayesian step on the agent's belief, projected back
onto a Gaussian family by maximum entropy
\citep{opper1996online,opper1998bayesian,caticha2020entropic}, and it carries
alongside the weights an explicit, dynamical estimate of how noisy each other
agent is as a channel --- a quantity we can read as trust
\citep{alves2016distrust,caticha2011agent}. Cognitive-sounding behaviour then
appears without being written in: surprise, blame attribution between
belief-revision and source-derision, and the ability to learn the \emph{opposite}
of what a distrusted source says.

Into this society we inject the smallest bias that can be correlated with a group
label. A fraction $\fd$ of agents shift their perceived agreement with a message
by $\pm d$ according to whether the sender shares their label. Nothing else
distinguishes the groups; the labels are informationally inert; the two groups are
treated symmetrically, so any asymmetry that emerges is spontaneous. The
resulting collective behaviour, as a function of $(d,\fd)$, is the subject of the
paper.

\paragraph{Contributions.}
\begin{itemize}\itemsep1pt
  \item \textbf{A phase diagram of emergent discrimination}
    (\S\ref{sec:results}): four regimes with sharp boundaries --- neutral
    polarisation unrelated to class, two discriminatory phases distinguished by
    whether ideology still tracks the group split, and a frustrated spin-glass
    regime when the bias favours the out-group. Discrimination needs a quorum:
    below $\fd\approx0.4$ the society stays neutral however strong the bias, and
    above it the transition in $d$ sharpens rapidly.
  \item \textbf{A mechanism specific to good learners}
    (\S\ref{sec:discrimination}): the bias displaces the separatrix of the
    learning flow, enlarging one dissonance-free basin at the other's expense.
    It acts only through the algorithm's trust sector, so optimising for the
    small-group setting is what makes a population vulnerable in the large-group
    one.
  \item \textbf{Which comes first, distrust or disagreement}
    (\S\ref{sec:results}): the order of polarisation is set by the agenda
    complexity $\alpha=P/K$, reversing at $\alpha\approx1.7$. This turns a
    contested question about polarisation
    \citep{fiorina2008polarization,iyengar2012affect} into a statement about a
    measurable parameter.
  \item \textbf{The same measurements in a society of language models}
    (\S\ref{sec:llm}). The order parameters need only pairwise opinion alignment
    and pairwise trust, both elicitable from language agents, so the phase
    structure is testable on a substrate the theory does not describe.
    \llmtodo{one sentence of findings once the study is run}
\end{itemize}

Our reading of the theoretical result is deliberately narrow: these are simple
machines, and extending conclusions from them to human societies is not
automatic. What the model does support is a claim about learning systems --- that
an inference rule optimal for a pair of agents, transplanted into a large
population, can convert an informationally worthless label into a stable
architecture of group distrust. For systems assembled from many interacting
learned agents that is a failure mode of the interaction rather than of the
training data, and it is invisible to any evaluation that examines agents one at a
time.
